\documentclass[book.tex]{subfiles}

\begin{document}

\chapter{Neural Tangent Kernel}
\label{chap:ntk}
\noindent\rule{11cm}{0.4pt} \\
\textbf{Prerequisites:}  \autoref{chap:ml_basics}, \autoref{chap:grad_descent}  \\
\textbf{Difficulty Level:} *** \\
\noindent\rule{11cm}{0.4pt}
\vspace{5mm}

The neural tangent kernel provides a theoretical methodology to study an "infinite-width" approximation of a fully connected network. Even more interestingly, the evolution of a neural network under gradient descent can also be described by a suitable kernel. 

\section{Introduction to Neural Tangent Kernels}

A recursion relation on neural networks can be mapped over to a relation on Gaussian processes. For a single layer network, we can write

\begin{align}
f_i^0(x) &= b_i^0 + \sum_j W^0_{ij} x_j \\
f_i^1(x) &= b_i^1 + \sum_j \sigma(W^0_{ij}) x_j \\
K^0(x, x') &= E[f_j(x)f_j(x')]
\end{align}

Let $f$ denote the output of the full network. For simplicity, assume that $f$ has a one dimensional output. Let $\theta_j$ denote the $j$-th learnable parameter. We then define the tangent kernel as

\begin{align}
\Theta(x, x') &= \sum_j \frac{\partial f(x)}{\partial \theta_j} \frac{\partial f(x')}{\partial \theta_j}
\end{align}

Note that the type of this kernel is
\begin{align}
\Theta: \mathbb{R}^n \times \mathbb{R}^n \to \mathbb{R}
\end{align}
where $x, y \in \mathbb{R}^n$. 

The neural tangent kernel lets us model the dynamics of the neural network during training. Let us define a dataset

\begin{align}
D &= \{(x_1, y_i),\dotsc, (x_N, y_N)\}
\end{align}

We can then model the time evolution of the output of the neural network $f$ by the function

\begin{align}
\frac{df(x; \theta)}{dt} &= - \sum_{i=1}^N \Theta(x, x_i) \frac{\partial \mathcal{L}(f(x_i; \theta), y_i)}{\partial f(x_i; \theta)} 
\end{align}

We can interpret the term $\Theta(x, x_i)$ as a measure of the influence of the gradient of the loss on the $i$-th datapoint $(x_i, y_i)$ on the output of the model. To first order, the evolution of the neural network during training g is given by the linear model
\begin{align}
f(x; \theta(t)) &= f(x; \theta(0)) + \sum_j \frac{\partial f(x; \theta_j(0))}{\partial \theta_j} (\theta_j(t) - \theta_j(0))
\end{align}

As the width of the neural network goes to infinity, its training behavior converges to this linear model. Note that for a finite width neural network, the neural tangent kernel $\Theta$ is a random quantity since it depends on the initialization of $\theta$. However for an infinitely wide network, the neural tangent kernel converges to a deterministic quantity. 
%What is the intuition behind why neural tangent kernel doesn't evolve? The gradient descent update equation vanishes as width $n \to \infty$. However gradient descent update for output is $O(1)$ so learning can still happen. So there's something interesting happening in the limiting behavior.

%The updates are small, the loss converges exponentially fast. Infinitely many model parameters each changing a vanishingly small amount have a collective effect.

%Here is a closed form for the evolution of the parameters.

%\begin{align}
%\theta_\mu(t) - \theta_\mu(0) &= - \frac{\partial f(\mathcal{X})}{\partial \theta_\mu} \Theta_0^{-1} \left ( I_{D \times D} - e^{-\Theta_0 t} \right ) \left ( f_0(\mathcal{X}) - \mathcal{y} \right)
%\end{align}

We can also compute the evolution of the parameters at a general point $x$. There is a Gaussian process that we can write down that describes the evolution of this system as a point process. This is called the NNGP kernel as opposed to the NNTK kernel.

%\begin{align}
%\\mu_t(x) &=  \Theta_0(x, \mathcal{X}) \Theta_0^{-1}\\
%\gamma_t(x) &= \\
%\end{align}

%Can do numerical comparisons with a wide resnet (\textcolor{red}{TODO}). Accuracy matches up to a point but you start to see some mismatches later as time dynamics evolve.

%(\textcolor{red}{TODO: Look at neural tangents library built on Jax to see how neural tangents work numerically. Might be a nice addition to DeepChem as well. This has a drop-in replacement for stax, one of the Jax nn libraries that allows for interesting drop-ins.})

When are NNTK/NNGP methods useful? They serve as a useful starting point for perturbation theoretic expansions. NNTK methods can also be competitive with neural network methods in their own right as predictive models.
\end{document}